% Proofs of Scope and Limitations
\section{Proofs of Scope and Limitations}

\subsection{Expressiveness Bounds}

\subsubsection{Theorem: Computational Power}
The Indus operational system has computational power equivalent to a deterministic finite automaton (DFA).

\textbf{Proof}:
\begin{enumerate}
    \item The system can be modeled as a DFA (construction shown in previous chapters)
    \item All attested inscriptions are recognized by this DFA
    \item No inscription requires greater computational power
    \item Therefore, the system is DFA-equivalent $\square$
\end{enumerate}

\subsection{Non-Expressibility Results}

\subsubsection{Theorem: No Context-Free Power}
The Indus system cannot express context-free languages.

\textbf{Proof}: Context-free languages require a pushdown automaton (PDA) with stack memory. The Indus system:
\begin{itemize}
    \item Has no nested structures (no stack)
    \item Shows no center-embedded patterns
    \item Exhibits no counting beyond finite bounds
    \item Therefore cannot be context-free $\square$
\end{itemize}

\subsubsection{Theorem: No Recursion}
The Indus system cannot express recursive structures.

\textbf{Proof}: Recursion requires:
\begin{itemize}
    \item Unbounded nesting depth, OR
    \item Self-referential rules
\end{itemize}

The Indus system exhibits:
\begin{itemize}
    \item Bounded inscription length
    \item Acyclic state graph (no cycles, no self-reference)
    \item Fixed maximum nesting (qualification depth)
    \item Therefore no recursion $\square$
\end{itemize}

\subsection{Positive Capabilities}

\subsubsection{Theorem: Regular Language Recognition}
The Indus system can recognize all patterns expressible as regular expressions.

\textbf{Proof}: By equivalence of DFA and regular expressions (Kleene's theorem), if the system is DFA-equivalent, it can recognize any regular language within its alphabet $\square$

\subsection{Cardinality Bounds}

\subsubsection{Theorem: Finite Language}
The set of all possible valid Indus inscriptions is finite.

\textbf{Proof}:
\begin{itemize}
    \item Alphabet $\Sigma$ is finite ($|\Sigma| \approx 417$)
    \item Maximum inscription length is bounded ($L_{\max} \approx 12$)
    \item Total possible inscriptions $\leq |\Sigma|^{L_{\max}}$
    \item This is finite $\square$
\end{itemize}

\subsection{Decidability Results}

\subsubsection{Theorem: Membership Decidability}
For any string $w$, we can decide in linear time whether $w$ is a valid inscription.

\textbf{Proof}: DFA recognition is $O(|w|)$. For each sign, perform one state transition. After $|w|$ steps, check if current state is accepting $\square$

\subsubsection{Theorem: Equivalence Decidability}
We can decide whether two inscriptions are equivalent under the operational model.

\textbf{Proof}: Two inscriptions are equivalent if they:
\begin{enumerate}
    \item Reach the same final state
    \item Perform the same sequence of operations
\end{enumerate}
Both properties are decidable by trace comparison $\square$

\subsection{Incompleteness Results}

\subsubsection{Theorem: Semantic Incompleteness}
The operational model cannot determine semantic content beyond functional categories.

\textbf{Statement}: This is not a proof but an observation: without bilingual texts or phonetic values, semantic interpretation is underdetermined by operational structure alone.

\subsection{Implications}

These formal results establish:
\begin{enumerate}
    \item The system's computational limits (DFA-level)
    \item What cannot be expressed (context-free structures, recursion)
    \item What can be effectively decided (validity, equivalence)
    \item The boundaries of formal analysis (operations, not semantics)
\end{enumerate}

These proofs validate the operational interpretation while acknowledging its inherent limitations.
